\section{Research Agenda}
\label{sec:research-agenda}

Revision 0.7 completes the first planned evidence stage, adds a structured
security capability matrix, and converts the nearest remaining work into a
public roadmap.

\subsection{Completed baseline: semantic mapping, fixed-input tests, and
parametric preservation}

The ordered SWAP semantics are now documented, all $(v,b,e,r_E)$ branches are
enumerated, the fifth-wire reduced states are calculated, and the victim
subsystem is compared with a clean routed reference. The resulting theorem is
\cref{eq:caseii-fixed-input-theorem}. Information metrics and deterministic
regression outputs are included in the public artifact.

Beyond the four BB84 inputs, victim-channel preservation is now numerically
verified across a 60-sample Bloch-sphere grid and a 200-sample pseudorandom
pure-state sweep, with cross-backend confirmation from the Qiskit Statevector
implementation. The use/read/retain/export privilege vocabulary has been
formalized and is used to label threat model cells.

\subsection{Immediate priority: exact detector reconstruction}

The strongest next scientific task is to identify the full educational
circuit's actual acceptance register. This requires:
\begin{enumerate}
  \item naming every classical selector, measurement result, and comparison bit;
  \item reconstructing the detector flag as a Boolean function;
  \item producing its truth table over every branch;
  \item replacing the provisional rule in
  \cref{eq:caseii-provisional-acceptance}; and
  \item regenerating the deterministic outputs and manuscript claims.
\end{enumerate}
This step should precede any strong statement that the original educational
detector has been bypassed.

\subsection{Completed: Qiskit parity and software decomposition}

The NumPy implementation is the validated reference. The Qiskit Statevector
backend now reproduces branch probabilities, reduced states, and metrics for all
tested parameterizations within documented tolerances (Issues \#2 and \#3
closed). A permanent \texttt{qiskit-parity} CI gate confirms this agreement on
every push. Complete branch-level parity across all circuit configurations
(Issue \#13) remains an open engineering task.

\subsection{Formal information-flow property}

After the detector semantics are fixed, define a noninterference-style property
for reclaimed qubits. For a protected variable $X$, attacker register $E$, and
approved victim observation policy $\mathcal{V}$, a secure transformation should
establish both victim equivalence and an explicit information bound such as
\begin{equation}
  \mathcal{V}(\rho_{\mathrm{attack}})
  \approx \mathcal{V}(\rho_{\mathrm{baseline}}),
  \qquad I(X:E)\leq\varepsilon.
\end{equation}
The present case study supplies a counterexample to histogram-only verification,
not yet a complete compiler-level noninterference framework. The matrix in
\cref{sec:capability-matrix} provides the initial domain of attacker access,
protected information, alternative routes, and mitigation annotations that the
formal property must eventually cover.

\subsection{Completed: malicious compiler prototype}

A controlled proof-of-concept compiler transformation has been implemented in
\texttt{src/quantum\_reuse/compiler\_pass.py}. The prototype inserts
\texttt{CNOT}($q_{\mathrm{signal}} \to q_{\mathrm{ancilla}}$) after Alice's
signal preparation and before any routing, demonstrating the \textbf{use}
$\to$ \textbf{retain} privilege conversion for both BB84 and arbitrary
pure-state inputs. Key results:
\begin{itemize}
  \item Victim preservation holds in 100\% of 260 tested inputs
    ($|\Delta F| < 10^{-10}$).
  \item $Z$-basis: injected ancilla = $\ket{v}\!\bra{v}$; attacker trace
    distance from $v{=}0$ to $v{=}1$ equals 1 (perfect distinguishability).
  \item $X$-basis: injected ancilla = $I/2$; attacker trace distance equals 0.
    Without Alice's basis, the single ancilla carries no value information.
\end{itemize}
The formal adversarial model bounding the joint \textbf{read}~+~\textbf{retain}
advantage over all single-ancilla injection strategies remains open
(Issue~23: ``Formalize optimal single-ancilla read-and-retain advantage'').
Issue~14 is closed for the constructive prototype; Issue~23 carries the open
theorem.

\subsection{Compiler prototype and hardware experiments}

Hardware and noise studies should compare coherent cleanup, measurement-reset
reuse, and fresh allocation, including routing, crosstalk, leakage, and
calibration drift. All experiments should use synthetic secrets and authorized
simulation or hardware.

\subsection{Venue and audience}

The work sits most naturally at the intersection of quantum software and quantum
security, with an educational circuit serving as the motivating case study. A
workshop, short research paper, or position paper aimed at that intersection is
a plausible first venue. Final venue selection should follow detector
reconstruction and external technical review.
